% ===================================================================
%  नकसा नं. १० — समुंदर। बंदरगाह।
%
%  Traduction, faite en 2026, en bhojpouri d'aujourd'hui, sur l'ido
%  de texto/io. Regles de la colonne : voir 10-nakasa-01.tex. Une
%  seule numerotation. 37 blocs, 98 renvois, vingt-neuf paires a
%  retourner.
%
%  LA RIVIERE EST A LUI, LA MER EST EMPRUNTEE, ET C'EST LA REPONSE
%  EXACTE DE LA COLONNE GUJARATIE. Elle avait releve, a ce meme
%  tableau 10, que son vocabulaire maritime etait ENTIEREMENT
%  gujarati — વહાણ le navire, સઢ la voile, ડોલકાઠી le mat, હલેસું la
%  rame, સુકાન le gouvernail — parce que le Gujarat est une cote et
%  que ses marchands tenaient Mombasa, Mascate et Malacca. Le pays
%  bhojpouri est l'inverse : huit cents kilometres de terre dans
%  toutes les directions, la langue majeure de l'Inde la plus
%  eloignee d'un rivage. Et le partage se lit mot par mot :
%
%    LE FLEUVE, EN PROPRE      LA MER, EMPRUNTEE
%    नाव    la barque          जहाज     le navire  (persan جهاز)
%    डोंगी  le canot           बंदरगाह  le port    (persan بندر)
%    पतवार  le gouvernail      लंगर     l'ancre    (persan لنگر)
%    डाँड़   la rame            मस्तूल   le mat     (arabe)
%    घाट    la berge           डेक, क्रेन, डॉक     (anglais)
%    माँझी  le batelier        महासागर  l'ocean    (sanskrit)
%    धार    le courant         ज्वार-भाटा la maree  (sanskrit)
%
%  LA LANGUE D'UN PEUPLE QUI NAVIGUE MAIS QUI NE PREND PAS LA MER.
%  Le Gange passe a la porte et tout ce qui le concerne se dit du
%  premier coup ; l'ocean commence exactement la ou commencent les
%  emprunts. Deux colonnes indiennes, deux tableaux 10, et la
%  geographie se lit dans le lexique sans qu'on ait besoin d'une
%  carte.
%
%  MAIS LE SOUS-MARIN (10) EST DE FABRICATION LOCALE, ET C'EST LE
%  SEUL. Sur les quarante-cinq objets de mer de ce tableau, un seul
%  mot moderne n'a ete emprunte a personne : « पनडुब्बी », de पानी
%  l'eau et डूबल couler — celui-qui-plonge-dans-l'eau. Il est fait
%  d'un verbe que la plaine avait deja, et c'est pour cela qu'elle a
%  pu le faire. On n'emprunte pas ce qu'on peut dire avec ce qu'on a.
%
%  Enfin l'arsenal (21) s'ecrit « सिलहखाना », le magasin d'armes, du
%  persan silāh-khāna. La colonne levantine avait releve au meme
%  endroit que « الترسانة » etait parti d'Italie et revenu en arabe ;
%  ici le mot n'a pas voyage par l'Europe du tout — il est venu par
%  la route de terre, avec les Moghols, et n'en est jamais reparti.
% ===================================================================

%%K t10-apar-1 apar
\VUsaut{4.0mm}
\VUtitre{15.0pt}{60}{नकसा नं. १०}
\VUsaut{3.0mm}
\VUpk{11.4pt}{40}{समुंदर। --- बंदरगाह।}
\VUsaut{3.0mm}

%%K t10-01-1 p
१. --- गरमी में, जब कुछ लोग पहाड़ पर सांति आ ठंढक खोजेला, दोसर लोग
समुंदर के किनारे जाइल पसंद करेला। पिछला साल हमनी के इहे कइनी जा,
काहेकि हमनी के \VUgras{महासागर} \textsuperscript{(1)} बहुते नीक लागेला।

%%K t10-02-1 p
२. --- हमरा त बहुते मजा आवेला ओह पानी के अपार फइलाव के निहारे में जवन
\VUgras{छितिज} \textsuperscript{(2)} तक फइलल बा, आ ओह बड़हन \VUgras{भाप
के जहाज} \textsuperscript{(3)} के लमहर \VUgras{समुंदरी जातरा} पर जात
देखे में, जवना पर अमेरिका जाए वाला जातरी चढ़ेलें।

%%K t10-03-1 p
३. --- एहीसे हमार बाबूजी, जे हमार मन जानेलें, छुट्टी बितावे खातिर
हरमेसा कवनो बहुते सांत बालू के किनारा चुनेलें, जवन हमनी के कवनो बड़का
\VUgras{फौजी बंदरगाह} के लगे होखे।

%%K t10-03-2 p
हमनी के पिछला बेर के ठहराव में \VUgras{जहाजी बेड़ा} ओह बड़हन
\VUgras{बाँध} \textsuperscript{(4)} के पाछे लंगर डालल, जवन \VUgras{फौजी
बंदरगाह} \textsuperscript{(5)} के घेरेला आ बचावेला, आ हम आपन मन भर के
तरह-तरह के \VUgras{लड़ाकू जहाज} देख पवनी — ऊ छोट \VUgras{पनडुब्बी}
\textsuperscript{(10)} जे डूब के पानी के नीचे गायब हो जाले, आ ऊ बड़हन
\VUgras{कवच वाला जहाज} \textsuperscript{(9)} जेकरा के अबहीं तक लगभग खाली
\VUgras{तारपीडो जहाज} \textsuperscript{(8)} के हमला के डर रहल।

%%K t10-tit-2 sub
\VUsaut{3.0mm}
\VUpk{10.2pt}{40}{छोट जहाज।}
\VUsaut{3.0mm}

%%K t10-04-1 p
४. --- ई छोट जहाज पहिलही से बहुते \VUgras{चतुराई} से मोटका धातु के कवच
के कमजोर जगहा खोज लेत रहे, ताकि ओहिजा खतरनाक \VUgras{तारपीडो} चिपका
सके, जेकर धमाका जहाज के डुबा देला ; तबो ऊ पनडुब्बी से कम डरावन रहे,
काहेकि तारपीडो-नासक जहाज ओकर पीछा आसानी से कर लेला।

%%K t10-05-1 p
५. --- हम कवच वाला तेज \VUgras{क्रूजर} \textsuperscript{(6)} भी देख पवनी,
जे लगभग ओतने अभेद्य बा जेतना असली कवच वाला जहाज, आ कई गो \VUgras{ढोवे
वाला जहाज} \textsuperscript{(12)}, तीन-चार \VUgras{तोप-नाव}
\textsuperscript{(7)} आ आखिर में एगो \VUgras{फ्रिगेट} \textsuperscript{(11)}।
\VUgras{एह बेड़ा के नजारा, जब ऊ लंगर उठावे खातिर चाल चलेला, सबसे
मजेदार होला।}

%%K t10-06-1 p
६. --- ओह बड़हन इस्पात के ढेर आ ओह सुंदर तीन मस्तूल वाला भा ओह लचकदार
\VUgras{पाल वाला जहाज} \textsuperscript{(13)} के बनावट में केतना अंतर बा,
जवन अबहीं ले बेपार के बेड़ा में चलत बा, आ जेकरा के हम पूरा पाल तानले
बंदरगाह में आवत देखनी जब हम \VUgras{घाट} \textsuperscript{(14)} पर टहलत
रहनी। सचहूँ ई जहाज आपन बहुते फायदा खो देले बाड़ें \VUgras{जब}
\VUgras{हवा} उलटा होखे। ओह लोग के सब पाल उतार के फेरु बंदरगाह में आवे
के पड़ल आ ओकनी के खींच के घाट तक ले आवे खातिर एगो \VUgras{खींचे वाला
जहाज} \textsuperscript{(16)} जरूरी भइल, जइसे ऊ सादा \VUgras{मालवाहक नाव}
\textsuperscript{(17)} होखसु।

%%K t10-tit-3 sub
\VUsaut{3.0mm}
\VUpk{10.2pt}{40}{बेपार के बंदरगाह।}
\VUsaut{3.0mm}

%%K t10-07-1 p
७. --- जवन बंदरगाह के बारे में हम कहत बानी ओहमें ई बात मजेदार बा कि
ओहमें खाली फौजी बंदरगाह नइखे, जवन बड़हन काम से बनावल गइल बा, बलुक एगो
अद्भुत \VUgras{बेपार के बंदरगाह} भी बा जवन एगो बड़का नदी के
\VUgras{मुहाना} में बा।

%%K t10-08-1 p
८. --- जवन जमीन के पट्टी दुनो बंदरगाह के अलग करेला ऊ किलाबंद बा आ
\VUgras{तोपखाना} आ \VUgras{बुर्ज} के कतार से बचावल बा, जवन समुंदर की ओर
बढ़ल बा। \VUgras{परकोटा} \textsuperscript{(15)} पर टहले खातिर खास इजाजत
चाहीं। हमरा ऊ आपन चाचा के जरिए आसानी से मिल गइल, जे \VUgras{जहाज बनावे
के कारखाना} \textsuperscript{(18)} के इंजीनियर हउवें, आ हम ओह जहाजन के
देखे गइनी जेकरा के बड़हन छप्पर के नीचे बनावल जात बा।

%%K t10-09-1 p
९. --- हम त एगो कवच वाला जहाज के पानी में उतारल भी देखनी, आ हमरा ऊ
भावुक पल इयाद बा जवन पूरा भीड़ महसूस कइलस जब ऊ बड़हन ढेर, आपन हाल पर
छोड़ल, आपन ढाँचा पर सरकल आ नदी के पानी पर तइरे लागल।

%%K t10-10-1 p
१०. --- कारखाना तक जाए खातिर \VUgras{कवायद के मैदान} \textsuperscript{(19)}
पार करे के पड़ेला, जहाँ छावनी के सिपाही रोज अभ्यास करे आवेला, आ एगो
लोहा के \VUgras{पुलिया} \textsuperscript{(20)} पर से गुजरे के पड़ेला, जवन
परेड के मैदान के \VUgras{सिलहखाना} \textsuperscript{(21)} से जोड़ेला।

%%K t10-tit-4 sub
\VUsaut{3.0mm}
\VUpk{10.2pt}{40}{सिलहखाना आ डॉक।}
\VUsaut{3.0mm}

%%K t10-11-1 p
११. --- आपन चाचा के साथे हम ओह बड़हन इमारत के देखनी जवन \VUgras{तोप}
\textsuperscript{(22)}, \VUgras{गोला} \textsuperscript{(23)} आ
\VUgras{गोला-बारूद} से भरल बा। हम ऊ \VUgras{ढलाईखाना}
\textsuperscript{(24)} भी देखनी जवना में भाप के जहाज के \VUgras{बॉयलर}
\textsuperscript{(25)} बनेला।

%%K t10-12-1 p
१२. --- एकदम लगे \VUgras{डॉक} \textsuperscript{(26)} बा — मरम्मत के डॉक
— आ ऊ बहुते देखे लायक \VUgras{तैरत डॉक} \textsuperscript{(27)}, जवना में
हम मरम्मत होखे वाला जहाज पर एकदम लगे से जहाज के \VUgras{पंखा}
\textsuperscript{(28)}, \VUgras{पेंदी} \textsuperscript{(29)} आ
\VUgras{पतवार} \textsuperscript{(30)} देख पवनी।

%%K t10-13-1 p
१३. --- बेपार के बंदरगाह भी ओतने देखे लायक बा जेतना फौजी बंदरगाह, आ हम
कई घंटा घाट पर ताकत बितवनी, जहाँ ऊ \VUgras{चाक वाला जहाज}
\textsuperscript{(31)} बान्हल बा जे नदी के उलटा चढ़ेला, आ ऊ
\VUgras{मस्तूल के क्रेन} \textsuperscript{(32)} के चलत देखनी, जे पंख नियर
बड़हन लदान उठा लेला।

%%K t10-tit-5 sub
\VUsaut{4.0mm}
\VUpk{10.2pt}{40}{रस्ता देखावे वाला।}
\VUsaut{3.0mm}

%%K t10-14-1 p
१४. --- हम ओह \VUgras{अटलांटिक जहाज} \textsuperscript{(34)} के देख के
दंग रह गइनी जे आपन \VUgras{झंडा} \textsuperscript{(35)} खुला हवा में
फहरावत बंदरगाह से निकलत रहे, आ जेकरा के एगो कुसल \VUgras{रस्ता देखावे
वाला मल्लाह} \textsuperscript{(36)} चलावत रहे, जे \VUgras{बोया}
\textsuperscript{(40)} आ निसान-खंभा से चिन्हल गहिर धार के साथे-साथे कमाल
के ढंग से चले जानत रहे, आ ओह \VUgras{चौड़ी नाव} \textsuperscript{(41)} से
बचत रहे जेकरा के आपन एकलौता चौकोर पाल से मुसकिल से चलावल जात रहे। हम
\VUgras{अगवारी} \textsuperscript{(37)} पर — आगे के हिस्सा — दुसरका दर्जा
के \VUgras{केबिन} \textsuperscript{(39)} पहिचननी, आ \VUgras{पिछवारी}
\textsuperscript{(38)} पर — पाछे के हिस्सा — पहिला दर्जा के जातरी लोग के
कमरा।

%%K t10-15-1 p
१५. --- कबो-कबो हम \VUgras{मालवाहक नाव} \textsuperscript{(42)} पर लदान
चढ़त भी देखनी, जवना में अभिए \VUgras{गोदाम} \textsuperscript{(43)} आ
\VUgras{समुंदरी इस्टेसन} \textsuperscript{(44)} के माल जमा कइल गइल रहे,
जवन \VUgras{घाट के पटरी} \textsuperscript{(45)} पर कई गो गाड़ी से आइल
रहे।

%%K t10-tit-6 sub
\VUsaut{3.0mm}
\VUpk{10.2pt}{40}{डाँड़ के मजा।}
\VUsaut{3.0mm}

%%K t10-16-1 p
१६. --- हम अकसर \VUgras{तैरे के डॉक} \textsuperscript{(53)} में नाव
चलवनी, जहाँ \VUgras{नाव} \textsuperscript{(46)} सरन लेवे आवेली, आ हमरा
\VUgras{डाँड़} \textsuperscript{(47)} चलावे में बहुते मजा आवत रहे। हम एगो
\VUgras{डेक} \textsuperscript{(48)} पर चढ़नी, जवन \VUgras{पाल वाला
नाव} \textsuperscript{(49)} के रहे, हम \VUgras{रस्सा} \textsuperscript{(52)}
के सहारे \VUgras{मस्तूल} \textsuperscript{(51)} पर \VUgras{आड़ा डंडा}
\textsuperscript{(50)} तक चढ़नी, आ ओकरा बाद हम आपन घुमाई \VUgras{डोंगी
से} \textsuperscript{(54)} फेरु कइनी, भारी \VUgras{मछरी मारे के नाव}
\textsuperscript{(55)} के बीच, जहाँ \VUgras{जाल} \textsuperscript{(56)}
सुखत रहे।

%%K t10-17-1 p
१७. --- \VUgras{घाट} \textsuperscript{(58)} पर हमरा सामने तरह-तरह के
\VUgras{माल} \textsuperscript{(59)} गुजरल, जवन \VUgras{पेटी}
\textsuperscript{(60)} में भा \VUgras{गठरी} \textsuperscript{(61)} में
रहे। ऊ बड़का नाव के \VUgras{लदान} \textsuperscript{(63)} बनावत रहे, आ
ताकतवर \VUgras{क्रेन} \textsuperscript{(62)} ओकरा के उठा के \VUgras{ट्रक}
\textsuperscript{(64)} पर धरत रहे, जबकि ट्रक वाला लोग ओकर घोसना करत आ
\VUgras{चुंगीघर} \textsuperscript{(65)} में चुंगी भरत रहे।

%%K t10-18-1 p
१८. --- आखिर में, जब हम \VUgras{घुमत पुल} \textsuperscript{(66)} भा
\VUgras{जल-फाटक} \textsuperscript{(69)} तक पहुँचनी, ओह \VUgras{गाद
निकाले के मसीन} \textsuperscript{(68)} के सामने से गुजरत जे \VUgras{नहर}
\textsuperscript{(67)} साफ करेले, हम \VUgras{इसारा-खंभा}
\textsuperscript{(70)} लगे घाट पर उतरनी।

%%K t10-19-1 p
१९. --- ओह घाट के दोसरा ओर ऊ \VUgras{बड़का डोंगी} \textsuperscript{(71)}
लागेली जे बेड़ा के अफसर लोग के किनारे ले आवेली। मल्लाह लोग के ताल से
आपन डाँड़ उठावत देखल देखे लायक नजारा बा, जबकि \VUgras{लहर}
\textsuperscript{(72)} पर डोलत नाव आसानी से उतरे के सीढ़ी लगे लाग जाले।
एह जगहा से \VUgras{ज्वार-भाटा} आ \VUgras{बालू के किनारा}
\textsuperscript{(73)} पर चढ़ाव आ उतार के चाल नीक से देखल जा सकेला।

%%K t10-tit-7 sub
\VUsaut{3.0mm}
\VUpk{10.2pt}{40}{कलब घर।}
\VUsaut{3.0mm}

%%K t10-20-1 p
२०. --- घर लवटे खातिर हम \VUgras{बिजली के ट्राम} \textsuperscript{(74)}
चलवनी, जेकर \VUgras{ठहराव} \textsuperscript{(75)} ओह \VUgras{खंभा} लगे बा
जवन अफसर लोग के कलब के सामने गाड़ल बा।

%%K t10-20-2 p
कलब के \VUgras{बालकनी} \textsuperscript{(76)} से, जवन \VUgras{लंगर}
\textsuperscript{(77)}, तोप आ गोला से सजल बा, हम अकसर जहाजी अफसर लोग के
सुंदर जमघट देखनी : \VUgras{लेफ्टिनेंट} \textsuperscript{(78)}, आपन
\VUgras{सीधा तलवार} \textsuperscript{(79)} के साथे, \VUgras{तोपखाना के
कप्तान} आ \VUgras{पैदल सेना} \textsuperscript{(80)} के कप्तान, आपन
\VUgras{बाँकी तलवार} \textsuperscript{(81)} के साथे। ओह लोग के पाछे एगो
\VUgras{मल्लाह} रहे जे ओह लोग के \VUgras{दूरबीन} देत रहे, जब ऊ लोग दूर
के कुछ देखल चाहत रहे।

%%K t10-21-1 p
२१. --- हम जवान \VUgras{उप-लेफ्टिनेंट} \textsuperscript{(82)} लोग के भी
देखनी जे आपन \VUgras{कमांडर} \textsuperscript{(83)} भा \VUgras{एडमिरल}
\textsuperscript{(84)} से हुकुम लेत रहे, जबकि एगो \VUgras{क्वार्टरमास्टर}
\textsuperscript{(85)} ओकरा के जहाज पर ले जाए खातिर अगोरत रहे।

%%K t10-22-1 p
२२. --- ओह बालकनी से नजारा शानदार बा : पूरा \VUgras{बालू के किनारा}
आपन \VUgras{तंबू} \textsuperscript{(86)} आ \VUgras{कपड़ा बदले के कोठरी}
\textsuperscript{(87)} के साथे नजर में आ जाला, आ नहाए के बेरा
\VUgras{नहाए वाला} \textsuperscript{(88)} लोग के \VUgras{जुआघर}
\textsuperscript{(89)} के सामने खुसी से उछलत देखल जा सकेला।

%%K t10-23-1 p
२३. --- किनारा के छोर पर तट एगो छोट \VUgras{चट्टानी प्रायद्वीप}
\textsuperscript{(91)} बनावेला, जहाँ \VUgras{बड़का लहर} \textsuperscript{(92)}
आके टूटेला आ जवना पर एगो \VUgras{दीपघर} \textsuperscript{(93)} बनल बा जे
रात में जहाजन के रस्ता देखावेला। ई प्रायद्वीप जमीन से खाली एगो बहुते
पातर \VUgras{जलडमरूमध्य} \textsuperscript{(90)} से जुड़ल बा।

%%K t10-tit-8 sub
\VUsaut{3.0mm}
\VUpk{10.2pt}{40}{तट के आम नजारा।}
\VUsaut{3.0mm}

%%K t10-24-1 p
२४. --- \VUgras{खड़ी चट्टान} \textsuperscript{(94)} दूर तक फइलल बा, जेकरा
के समुंदर चीर देले बा आ ओहमें मनमाना ढंग से \VUgras{खाड़ी}
\textsuperscript{(95)} आ \VUgras{अंतरीप} के कतार बना देले बा, जवना से ऊ
बहुते सुंदर लागेला।

%%K t10-24-2 p
\VUgras{पठार} \textsuperscript{(96)} पर, जवन \VUgras{जलसंधि}
\textsuperscript{(98)} के ऊपर बा, एगो बहुते जरूरी \VUgras{किला}
\textsuperscript{(97)} आ कुछ « कोठी » बा।

%%K t10-25-1 p
२५. --- ई जलसंधि महादीप से एगो बहुते खतरनाक \VUgras{टापू}
\textsuperscript{(99)} के अलग करेला, जेकरा चारो ओर \VUgras{चट्टान के
रोड़ा} \textsuperscript{(100)} आ लहर तोड़त चट्टान बा, जहाँ नाव आ बड़का
जहाज तक कबो-कबो फँस जाला। मदद के तइयारी आ \VUgras{बचाव के नाव} के
जल्दी पहुँचला के बादो ई \VUgras{जहाज के डूबल} डरावन होला, आ बराबरी के
दिन के \VUgras{तूफान} में कई गो जहाज मनई आ माल के साथे डूब गइल बा।

%%K t10-25-2 p
एह लगे होखला के बादो \VUgras{बंदरगाह} में मल्लाह लोग के बहुते आवाजाही
रहेला।
